\section{Conclusion}
\label{sec:conclusion}

LLMBO-MO incorporates battery-domain knowledge into multiobjective charging
optimization through initial protocol proposals and regional preferences for
successive scalarized subproblems. The Region-Lift adaptation uses GP
posterior cross-covariance to translate the regional preferences into mean
adjustments when computing EI.

With 56 battery-model evaluations, LLMBO-MO achieved a mean final HV of
\ChenLLMBOMean\ on Chen2020, compared with \ChenParEGOMean\ for ParEGO,
a \ChenFinalGain\ increase. The corresponding Ecker2015 values were
\EckerLLMBOMean\ and \EckerParEGOMean. In the 16-evaluation component study,
initialization accounted for most of the improvement, while adding regional
guidance increased the mean HV from \AblationWarmMean\ to \AblationFullMean.
The initialization study further showed that battery-specific prompting
improved short-budget optimization more than generic prompting. These findings
support the use of qualitative battery knowledge to improve the early search
for charging trade-offs. Future work will assess regional guidance over wider
operating conditions and validate the resulting protocols on physical cells
with measured degradation.
